\chapter{Casos de uso adicionales}\label{cap:casos_de_uso_adicionales}

Este apéndice recoge los casos de uso adicionales del sistema que, por no considerarse esenciales para la comprensión general del proyecto, no se han incluido en el cuerpo principal de la memoria (Sección~\ref{cap:casos_de_uso}). Cada uno de ellos sigue la misma estructura que los presentados en dicha sección: resumen, actor, precondición, postcondición, secuencia base, secuencia alternativa y requisitos relacionados.

\begin{table}[h!]
    \centering
    \renewcommand{\arraystretch}{1.5}

    \begin{tabular}{|p{4cm}|p{11cm}|}
        \hline
        \multicolumn{2}{|l|}{\textbf{Caso de uso: Consultar información de una cofradía, procesión, evento o paso}}                                                                              \\
        \hline
        \textbf{Resumen}                                                                                                                                           & El usuario consulta la información detallada de una cofradía, procesión, evento o paso de la ciudad seleccionada.                               \\
        \hline
        \textbf{Actor}                                                                                                                                             & Ciudadano                                                                                                              \\
        \hline
        \textbf{Precondición}                                                                                                                                      & El usuario ha seleccionado una ciudad o el usuario ha autorizado que la aplicación pueda consultar su ubicación actual. \\
        \hline
        \textbf{Postcondición}                                                                                                                                     & Se muestra la lista de cofradías con su información.                                                                   \\
        \hline
        \multicolumn{2}{|l|}{\textbf{Secuencia Base}}                                                                                                               \\
        \hline
        \multicolumn{2}{|p{15cm}|}{ [1] El usuario accede al listado de cofradías, procesiones, eventos o pasos (RF-37). }                                                                              \\
        \hline
        \multicolumn{2}{|p{15cm}|}{ [2] El sistema muestra la lista correspondiente de la ciudad seleccionada. }                                                       \\
        \hline
        \multicolumn{2}{|p{15cm}|}{ [3] El usuario selecciona un elemento de la lista. }                                                                                       \\
        \hline
        \multicolumn{2}{|p{15cm}|}{ [4] El sistema muestra la información detallada del elemento seleccionado. }                                                               \\
        \hline
        \multicolumn{2}{|l|}{\textbf{Secuencia Alternativa}}                                                                                                        \\
        \hline
        \multicolumn{2}{|p{15cm}|}{ [2.a.1] Sin elementos disponibles: El sistema muestra un mensaje indicando que no hay elementos registrados de ese tipo en esa ciudad. } \\
        \hline
        \textbf{Requisitos Relacionados}                                                                                                                           & RF-03, RF-04, RF-05 y RF-17.                                                                                                  \\
        \hline
    \end{tabular}
    \caption{Caso de uso: Consultar información de una cofradía, procesión, evento o paso}
\end{table}


\begin{table}[h!]
    \centering
    \renewcommand{\arraystretch}{1.5}

    \begin{tabular}{|p{4cm}|p{11cm}|}
        \hline
        \multicolumn{2}{|l|}{\textbf{Caso de uso: Acceso como cofrade}}                                                                   \\
        \hline
        \textbf{Resumen}                                                                                                                 & Un usuario se registra en la aplicación como miembro de una cofradía. \\
        \hline
        \textbf{Actor}                                                                                                                   & Cofrade                                                               \\
        \hline
        \textbf{Precondición}                                                                                                            & El usuario no está registrado.                                        \\
        \hline
        \textbf{Postcondición}                                                                                                           & El usuario queda registrado y puede acceder a funciones de cofrade.   \\
        \hline
        \multicolumn{2}{|l|}{\textbf{Secuencia Base}}                                                                                     \\
        \hline
        \multicolumn{2}{|p{15cm}|}{ [1] El usuario selecciona la opción de registro. }                                                    \\
        \hline
        \multicolumn{2}{|p{15cm}|}{ [2] El sistema muestra el formulario de registro (Introducir código de acceso para dicha cofradía). } \\
        \hline
        \multicolumn{2}{|p{15cm}|}{ [3] El usuario introduce el código de acceso de la cofradía. }                                        \\
        \hline
        \multicolumn{2}{|p{15cm}|}{ [4] El sistema valida el código. }                                                                    \\
        \hline
        \multicolumn{2}{|p{15cm}|}{ [5] El sistema muestra mensaje informando del correcto acceso. }                                      \\
        \hline
        \multicolumn{2}{|l|}{\textbf{Secuencia Alternativa}}                                                                              \\
        \hline
        \multicolumn{2}{|p{15cm}|}{ [3.a.1] Cancelación de la operación: El usuario desea cancelar la operación. }                           \\
        \hline
        \multicolumn{2}{|p{15cm}|}{ [5.a.1] Código inválido: El sistema muestra los errores de validación. }                              \\
        \hline
        \textbf{Requisitos Relacionados}                                                                                                 & RF-02.                                                                \\
        \hline
    \end{tabular}
    \caption{Caso de uso: Acceso como cofrade}
\end{table}


\begin{table}[h!]
    \centering
    \renewcommand{\arraystretch}{1.5}

    \begin{tabular}{|p{4cm}|p{11cm}|}
        \hline
        \multicolumn{2}{|l|}{\textbf{Caso de uso: Añadir favoritos}}                                                                    \\
        \hline
        \textbf{Resumen}                                                                                                               & Un usuario marca elementos como favoritos. \\
        \hline
        \textbf{Actor}                                                                                                                 & Ciudadano                                  \\
        \hline
        \textbf{Precondición}                                                                                                          &                                            \\
        \hline
        \textbf{Postcondición}                                                                                                         & El elemento queda guardado de favoritos    \\
        \hline
        \multicolumn{2}{|l|}{\textbf{Secuencia Base}}                                                                                   \\
        \hline
        \multicolumn{2}{|p{15cm}|}{ [1] El usuario navega a una cofradía, evento o procesión. }                                         \\
        \hline
        \multicolumn{2}{|p{15cm}|}{ [2] El usuario pulsa el botón de favorito. }                                                        \\
        \hline
        \multicolumn{2}{|p{15cm}|}{ [3] El sistema guarda el elemento como favorito. }                                                  \\
        \hline
        \multicolumn{2}{|p{15cm}|}{ [4] El sistema muestra confirmación visual. }                                                       \\
        \hline
        \multicolumn{2}{|l|}{\textbf{Secuencia Alternativa}}                                                                            \\
        \hline
        \multicolumn{2}{|p{15cm}|}{ [2.a.1] El elemento ya es favorito: El sistema lo elimina de favoritos (caso de uso Eliminar favoritos). } \\
        \hline
        \textbf{Requisitos Relacionados}                                                                                               & RF-12.                                     \\
        \hline
    \end{tabular}
    \caption{Caso de uso: Añadir favoritos}
\end{table}


\begin{table}[h!]
    \centering
    \renewcommand{\arraystretch}{1.5}

    \begin{tabular}{|p{4cm}|p{11cm}|}
        \hline
        \multicolumn{2}{|l|}{\textbf{Caso de uso: Eliminar favoritos}}                                                                  \\
        \hline
        \textbf{Resumen}                                                                                                               & Un usuario desmarca elementos como favoritos. \\
        \hline
        \textbf{Actor}                                                                                                                 & Ciudadano                                      \\
        \hline
        \textbf{Precondición}                                                                                                          & El elemento esté marcado como favorito        \\
        \hline
        \textbf{Postcondición}                                                                                                         & El elemento queda eliminado de favoritos      \\
        \hline
        \multicolumn{2}{|l|}{\textbf{Secuencia Base}}                                                                                   \\
        \hline
        \multicolumn{2}{|p{15cm}|}{ [1] El usuario navega a una cofradía, evento o procesión. }                                         \\
        \hline
        \multicolumn{2}{|p{15cm}|}{ [2] El usuario pulsa el botón de favorito. }                                                        \\
        \hline
        \multicolumn{2}{|p{15cm}|}{ [3] El sistema quita el elemento de favorito. }                                                     \\
        \hline
        \multicolumn{2}{|p{15cm}|}{ [4] El sistema muestra confirmación visual. }                                                       \\
        \hline
        \multicolumn{2}{|l|}{\textbf{Secuencia Alternativa}}                                                                            \\
        \hline
        \multicolumn{2}{|p{15cm}|}{ [2.a.1] El elemento no es favorito: El sistema lo añade a favoritos (caso de uso Añadir favoritos). } \\
        \hline
        \textbf{Requisitos Relacionados}                                                                                               & RF-12.                                        \\
        \hline
    \end{tabular}
    \caption{Caso de uso: Eliminar favoritos}
\end{table}


\begin{table}[h!]
    \centering
    \renewcommand{\arraystretch}{1.5}

    \begin{tabular}{|p{4cm}|p{11cm}|}
        \hline
        \multicolumn{2}{|l|}{\textbf{Caso de uso: Configurar notificaciones}}                                   \\
        \hline
        \textbf{Resumen}                                                                                       & Un ciudadano se suscribe o desuscribe de las notificaciones sobre cofradías, eventos o procesiones concretas. \\
        \hline
        \textbf{Actor}                                                                                         & Ciudadano                                     \\
        \hline
        \textbf{Precondición}                                                                                  & --                                             \\
        \hline
        \textbf{Postcondición}                                                                                 & Las preferencias de notificaciones quedan guardadas \\
        \hline
        \multicolumn{2}{|l|}{\textbf{Secuencia Base}}                                                           \\
        \hline
        \multicolumn{2}{|p{15cm}|}{ [1] El usuario accede a la configuración de notificaciones. }               \\
        \hline
        \multicolumn{2}{|p{15cm}|}{ [2] El sistema muestra las opciones disponibles. }                          \\
        \hline
        \multicolumn{2}{|p{15cm}|}{ [3] El usuario activa/desactiva tipos de notificaciones. }                  \\
        \hline
        \multicolumn{2}{|p{15cm}|}{ [4] El usuario guarda los cambios. }                                        \\
        \hline
        \multicolumn{2}{|p{15cm}|}{ [5] El sistema confirma que se han guardado. }                              \\
        \hline
        \multicolumn{2}{|l|}{\textbf{Secuencia Alternativa}}                                                    \\
        \hline
        \multicolumn{2}{|p{15cm}|}{ [4.a.1] Cancelación de la operación: El usuario desea cancelar la operación. } \\
        \hline
        \textbf{Requisitos Relacionados}                                                                       & RF-22.                                        \\
        \hline
    \end{tabular}
    \caption{Caso de uso: Configurar notificaciones}
\end{table}


\begin{table}[h!]
    \centering
    \renewcommand{\arraystretch}{1.5}

    \begin{tabular}{|p{4cm}|p{11cm}|}
        \hline
        \multicolumn{2}{|l|}{\textbf{Caso de uso: Crear cofradía/Hermandad}}                                    \\
        \hline
        \textbf{Resumen}                                                                                       & Un miembro de la Junta de Cofradías crea una nueva cofradía. \\
        \hline
        \textbf{Actor}                                                                                         & Junta de Cofradías                                           \\
        \hline
        \textbf{Precondición}                                                                                  & El usuario tiene rol de Junta de Cofradías                   \\
        \hline
        \textbf{Postcondición}                                                                                 & La cofradía queda creada en el sistema                       \\
        \hline
        \multicolumn{2}{|l|}{\textbf{Secuencia Base}}                                                           \\
        \hline
        \multicolumn{2}{|p{15cm}|}{ [1] El usuario accede al panel de gestión. }                                \\
        \hline
        \multicolumn{2}{|p{15cm}|}{ [2] El usuario selecciona crear nueva cofradía. }                           \\
        \hline
        \multicolumn{2}{|p{15cm}|}{ [3] El sistema muestra el formulario de creación. }                         \\
        \hline
        \multicolumn{2}{|p{15cm}|}{ [4] El usuario introduce los datos de la cofradía. }                       \\
        \hline
        \multicolumn{2}{|p{15cm}|}{ [5] El usuario guarda la cofradía. }                                       \\
        \hline
        \multicolumn{2}{|p{15cm}|}{ [6] El sistema valida y guarda la información. }                            \\
        \hline
        \multicolumn{2}{|l|}{\textbf{Secuencia Alternativa}}                                                    \\
        \hline
        \multicolumn{2}{|p{15cm}|}{ [4.a.1] Cancelación de la operación: El usuario desea cancelar la operación. } \\
        \hline
        \multicolumn{2}{|p{15cm}|}{ [6.a.1] Datos incompletos: El sistema muestra los campos requeridos. }      \\
        \hline
        \textbf{Requisitos Relacionados}                                                                       & RF-21 y RF-33.                                               \\
        \hline
    \end{tabular}
    \caption{Caso de uso: Crear cofradía/Hermandad}
\end{table}


\begin{table}[h!]
    \centering
    \renewcommand{\arraystretch}{1.5}

    \begin{tabular}{|p{4cm}|p{11cm}|}
        \hline
        \multicolumn{2}{|l|}{\textbf{Caso de uso: Crear procesión}}                                             \\
        \hline
        \textbf{Resumen}                                                                                       & Un miembro de la Junta de Cofradías crea una nueva procesión.                  \\
        \hline
        \textbf{Actor}                                                                                         & Junta de Cofradías                                                             \\
        \hline
        \textbf{Precondición}                                                                                  & El usuario tiene rol de Junta de Cofradías y que exista al menos una cofradía \\
        \hline
        \textbf{Postcondición}                                                                                 & La procesión queda creada en el sistema                                        \\
        \hline
        \multicolumn{2}{|l|}{\textbf{Secuencia Base}}                                                           \\
        \hline
        \multicolumn{2}{|p{15cm}|}{ [1] El usuario accede al panel de gestión. }                                \\
        \hline
        \multicolumn{2}{|p{15cm}|}{ [2] El usuario selecciona crear nueva procesión. }                          \\
        \hline
        \multicolumn{2}{|p{15cm}|}{ [3] El sistema muestra el formulario de creación. }                         \\
        \hline
        \multicolumn{2}{|p{15cm}|}{ [4] El usuario introduce los datos de la procesión. }                       \\
        \hline
        \multicolumn{2}{|p{15cm}|}{ [5] El usuario define el recorrido en el mapa. }                            \\
        \hline
        \multicolumn{2}{|p{15cm}|}{ [6] El usuario guarda la procesión. }                                       \\
        \hline
        \multicolumn{2}{|p{15cm}|}{ [7] El sistema valida y guarda la información. }                            \\
        \hline
        \multicolumn{2}{|l|}{\textbf{Secuencia Alternativa}}                                                    \\
        \hline
        \multicolumn{2}{|p{15cm}|}{ [4.a.1] Cancelación de la operación: El usuario desea cancelar la operación. } \\
        \hline
        \multicolumn{2}{|p{15cm}|}{ [5.a.1] Cancelación de la operación: El usuario desea cancelar la operación. } \\
        \hline
        \multicolumn{2}{|p{15cm}|}{ [7.a.1] Datos incompletos: El sistema muestra los campos requeridos. }      \\
        \hline
        \textbf{Requisitos Relacionados}                                                                       & RF-21 y RF-30.                                                                 \\
        \hline
    \end{tabular}
    \caption{Caso de uso: Crear procesión}
\end{table}


\begin{table}[h!]
    \centering
    \renewcommand{\arraystretch}{1.5}

    \begin{tabular}{|p{4cm}|p{11cm}|}
        \hline
        \multicolumn{2}{|l|}{\textbf{Caso de uso: Actualizar estado de procesión}}                                                 \\
        \hline
        \textbf{Resumen}                                                                                                          & La Junta de Cofradías actualiza el estado de una procesión.                             \\
        \hline
        \textbf{Actor}                                                                                                            & Junta de Cofradías                                                                      \\
        \hline
        \textbf{Precondición}                                                                                                     & El usuario tiene rol de Junta de Cofradías y existe una procesión programada o en curso \\
        \hline
        \textbf{Postcondición}                                                                                                    & El estado de la procesión queda actualizado y se notifica a los usuarios                \\
        \hline
        \multicolumn{2}{|l|}{\textbf{Secuencia Base}}                                                                              \\
        \hline
        \multicolumn{2}{|p{15cm}|}{ [1] El usuario accede al panel de gestión de procesiones. }                                    \\
        \hline
        \multicolumn{2}{|p{15cm}|}{ [2] El usuario selecciona la procesión a actualizar. }                                         \\
        \hline
        \multicolumn{2}{|p{15cm}|}{ [3] El sistema muestra la información actual de la procesión. }                                \\
        \hline
        \multicolumn{2}{|p{15cm}|}{ [4] El usuario selecciona el nuevo estado. }                                                   \\
        \hline
        \multicolumn{2}{|p{15cm}|}{ [5] El usuario confirma el cambio. }                                                           \\
        \hline
        \multicolumn{2}{|p{15cm}|}{ [6] El sistema actualiza el estado. }                                                          \\
        \hline
        \multicolumn{2}{|p{15cm}|}{ [7] El sistema envía notificaciones automáticas a los usuarios interesados. }                  \\
        \hline
        \multicolumn{2}{|l|}{\textbf{Secuencia Alternativa}}                                                                       \\
        \hline
        \multicolumn{2}{|p{15cm}|}{ [2.a.1] Sin conexión: El sistema muestra error de conexión. }                                  \\
        \hline
        \multicolumn{2}{|p{15cm}|}{ [2.b.1] Cancelación de la operación: El usuario desea cancelar la operación. }                    \\
        \hline
        \multicolumn{2}{|p{15cm}|}{ [4.a.1] Sin conexión: El sistema muestra error de conexión. }                                  \\
        \hline
        \multicolumn{2}{|p{15cm}|}{ [4.b.1] Caso de uso crear notificación: El usuario crea una notificación. }                                \\
        \hline
        \multicolumn{2}{|p{15cm}|}{ [4.c.1] Estado no válido: El sistema muestra los estados disponibles según el estado actual. } \\
        \hline
        \multicolumn{2}{|p{15cm}|}{ [4.d.1] Cancelación de la operación: El usuario desea cancelar la operación. }                    \\
        \hline
        \multicolumn{2}{|p{15cm}|}{ [5.a.1] Sin conexión: El sistema muestra error de conexión. }                                  \\
        \hline
        \multicolumn{2}{|p{15cm}|}{ [5.b.1] Cancelación de la operación: El usuario desea cancelar la operación. }                    \\
        \hline
        \textbf{Requisitos Relacionados}                                                                                          & RF-28 y RF-27.                                                                   \\
        \hline
    \end{tabular}
    \caption{Caso de uso: Actualizar estado de procesión}
\end{table}


\begin{table}[h!]
    \centering
    \renewcommand{\arraystretch}{1.5}

    \begin{tabular}{|p{4cm}|p{11cm}|}
        \hline
        \multicolumn{2}{|l|}{\textbf{Caso de uso: Gestionar información de Semana Santa}}                                       \\
        \hline
        \textbf{Resumen}                                                                                                       & La Junta de Cofradías introduce o modifica toda la información de la Semana Santa de su ciudad. \\
        \hline
        \textbf{Actor}                                                                                                         & Junta de Cofradías                                                                              \\
        \hline
        \textbf{Precondición}                                                                                                  & El usuario tiene rol de Junta de Cofradías                                                      \\
        \hline
        \textbf{Postcondición}                                                                                                 & La información de la Semana Santa queda actualizada en el sistema                               \\
        \hline
        \multicolumn{2}{|l|}{\textbf{Secuencia Base}}                                                                           \\
        \hline
        \multicolumn{2}{|p{15cm}|}{ [1] El usuario selecciona el módulo a gestionar (cofradías, procesiones, eventos, pasos). } \\
        \hline
        \multicolumn{2}{|p{15cm}|}{ [2] El sistema muestra la lista de elementos existentes. }                                  \\
        \hline
        \multicolumn{2}{|p{15cm}|}{ [3] El usuario puede crear, editar o eliminar elementos. }                                  \\
        \hline
        \multicolumn{2}{|p{15cm}|}{ [4] El usuario introduce/modifica los datos requeridos. }                                   \\
        \hline
        \multicolumn{2}{|p{15cm}|}{ [5] El sistema valida la información. }                                                     \\
        \hline
        \multicolumn{2}{|p{15cm}|}{ [6] El sistema guarda los cambios. }                                                        \\
        \hline
        \multicolumn{2}{|p{15cm}|}{ [7] El sistema muestra confirmación. }                                                      \\
        \hline
        \multicolumn{2}{|l|}{\textbf{Secuencia Alternativa}}                                                                    \\
        \hline
        \multicolumn{2}{|p{15cm}|}{ [3.a.1] Sin conexión: El sistema muestra error de conexión}                                 \\
        \hline
        \multicolumn{2}{|p{15cm}|}{ [3.b.1] Cancelación de la operación: El usuario desea cancelar la operación. }                 \\
        \hline
        \multicolumn{2}{|p{15cm}|}{ [4.a.1] Sin conexión: El sistema muestra error de conexión. }                               \\
        \hline
        \multicolumn{2}{|p{15cm}|}{ [4.b.1] Cancelación de la operación: El usuario desea cancelar la operación. }                 \\
        \hline
        \multicolumn{2}{|p{15cm}|}{ [4.c.1] Eliminar elemento con dependencias: El sistema advierte y solicita confirmación. }  \\
        \hline
        \multicolumn{2}{|p{15cm}|}{ [6.a.1] Datos inválidos: El sistema muestra errores de validación. }                        \\
        \hline
        \textbf{Requisitos Relacionados}                                                                                       & RF-20, RF-21 y RF-31.                                                                           \\
        \hline
    \end{tabular}
    \caption{Caso de uso: Gestionar información de Semana Santa}
\end{table}


\begin{table}[h!]
    \centering
    \renewcommand{\arraystretch}{1.5}

    \begin{tabular}{|p{4cm}|p{11cm}|}
        \hline
        \multicolumn{2}{|l|}{\textbf{Caso de uso: Definir recorrido de procesión}}                                                 \\
        \hline
        \textbf{Resumen}                                                                                                          & La Junta de Cofradías define o modifica el recorrido de una procesión en el mapa. \\
        \hline
        \textbf{Actor}                                                                                                            & Junta de Cofradías                                                                \\
        \hline
        \textbf{Precondición}                                                                                                     & El usuario tiene rol de Junta de Cofradías y existe una procesión creada.         \\
        \hline
        \textbf{Postcondición}                                                                                                    & El recorrido de la procesión queda definido con todos sus puntos.                 \\
        \hline
        \multicolumn{2}{|l|}{\textbf{Secuencia Base}}                                                                              \\
        \hline
        \multicolumn{2}{|p{15cm}|}{ [1] El usuario accede a la gestión de procesiones. }                                           \\
        \hline
        \multicolumn{2}{|p{15cm}|}{ [2] El usuario selecciona la procesión. }                                                      \\
        \hline
        \multicolumn{2}{|p{15cm}|}{ [3] El usuario accede a la edición del recorrido. }                                            \\
        \hline
        \multicolumn{2}{|p{15cm}|}{ [4] El sistema muestra el mapa interactivo. }                                                  \\
        \hline
        \multicolumn{2}{|p{15cm}|}{ [5] El usuario marca los puntos del recorrido en orden. }                                      \\
        \hline
        \multicolumn{2}{|p{15cm}|}{ [6] El usuario define horarios estimados para cada puntos. }                                   \\
        \hline
        \multicolumn{2}{|p{15cm}|}{ [7] El usuario marca puntos de interés especiales. }                                           \\
        \hline
        \multicolumn{2}{|p{15cm}|}{ [8] El usuario guarda el recorrido. }                                                          \\
        \hline
        \multicolumn{2}{|p{15cm}|}{ [9] El sistema calcula distancia y tiempo estimado. }                                          \\
        \hline
        \multicolumn{2}{|l|}{\textbf{Secuencia Alternativa}}                                                                       \\
        \hline
        \multicolumn{2}{|p{15cm}|}{ [2.a.1] Cancelación de la operación: El usuario desea cancelar la operación. }                    \\
        \hline
        \multicolumn{2}{|p{15cm}|}{ [3.a.1] Cancelación de la operación: El usuario desea cancelar la operación. }                    \\
        \hline
        \multicolumn{2}{|p{15cm}|}{ [5.a.1] Cancelación de la operación: El usuario desea cancelar la operación. }                    \\
        \hline
        \multicolumn{2}{|p{15cm}|}{ [5.a.2] Recorrido inválido: El sistema indica que debe tener al menos punto de inicio y fin. } \\
        \hline
        \multicolumn{2}{|p{15cm}|}{ [6.a.1] Cancelación de la operación: El usuario desea cancelar la operación. }                    \\
        \hline
        \multicolumn{2}{|p{15cm}|}{ [7.a.1] Cancelación de la operación: El usuario desea cancelar la operación. }                    \\
        \hline
        \multicolumn{2}{|p{15cm}|}{ [8.a.1] Cancelación de la operación: El usuario desea cancelar la operación. }                    \\
        \hline
        \textbf{Requisitos Relacionados}                                                                                          & RF-30 y RF-29.                                                                    \\
        \hline
    \end{tabular}
    \caption{Caso de uso: Definir recorrido de procesión}
\end{table}


\begin{table}[h!]
    \centering
    \renewcommand{\arraystretch}{1.5}

    \begin{tabular}{|p{4cm}|p{11cm}|}
        \hline
        \multicolumn{2}{|l|}{\textbf{Caso de uso: Dar de alta Junta de Cofradías}}                                                  \\
        \hline
        \textbf{Resumen}                                                                                                            & El administrador da de alta la Junta de Cofradías encargada de gestionar una ciudad.                                 \\
        \hline
        \textbf{Actor}                                                                                                              & Administrador                                                                                                        \\
        \hline
        \textbf{Precondición}                                                                                                       & El usuario tiene rol de Administrador y existe la ciudad a la que se va a asociar la Junta de Cofradías.             \\
        \hline
        \textbf{Postcondición}                                                                                                      & La nueva Junta de Cofradías queda creada y asociada a la ciudad seleccionada.                                        \\
        \hline
        \multicolumn{2}{|l|}{\textbf{Secuencia Base}}                                                                              \\
        \hline
        \multicolumn{2}{|p{15cm}|}{ [1] El usuario accede al panel de administración. }                                            \\
        \hline
        \multicolumn{2}{|p{15cm}|}{ [2] El sistema muestra la lista de ciudades existentes. }                                      \\
        \hline
        \multicolumn{2}{|p{15cm}|}{ [3] El usuario selecciona la ciudad para la que desea crear la Junta de Cofradías. }           \\
        \hline
        \multicolumn{2}{|p{15cm}|}{ [4] El usuario introduce los datos de la Junta de Cofradías. }                                 \\
        \hline
        \multicolumn{2}{|p{15cm}|}{ [5] El sistema valida y guarda la nueva Junta de Cofradías, asociándola a la ciudad seleccionada. } \\
        \hline
        \multicolumn{2}{|l|}{\textbf{Secuencia Alternativa}}                                                                       \\
        \hline
        \multicolumn{2}{|p{15cm}|}{ [3.a.1] Ciudad con Junta ya asignada: El sistema informa de que la ciudad ya tiene una Junta de Cofradías asociada. } \\
        \hline
        \multicolumn{2}{|p{15cm}|}{ [4.a.1] Cancelación de la operación: El usuario desea cancelar la operación. }                 \\
        \hline
        \multicolumn{2}{|p{15cm}|}{ [5.a.1] Datos incompletos: El sistema muestra los campos requeridos. }                         \\
        \hline
        \textbf{Requisitos Relacionados}                                                                                           & RF-34.                                                                                                                \\
        \hline
    \end{tabular}
    \caption{Caso de uso: Dar de alta Junta de Cofradías}
\end{table}


\begin{table}[h!]
    \centering
    \renewcommand{\arraystretch}{1.5}

    \begin{tabular}{|p{4cm}|p{11cm}|}
        \hline
        \multicolumn{2}{|l|}{\textbf{Caso de uso: Editar Junta de Cofradías}}                                                       \\
        \hline
        \textbf{Resumen}                                                                                                            & El administrador modifica los datos de una Junta de Cofradías existente.                                             \\
        \hline
        \textbf{Actor}                                                                                                              & Administrador                                                                                                        \\
        \hline
        \textbf{Precondición}                                                                                                       & El usuario tiene rol de Administrador y existe al menos una Junta de Cofradías registrada.                           \\
        \hline
        \textbf{Postcondición}                                                                                                      & Los datos de la Junta de Cofradías quedan actualizados en el sistema.                                                \\
        \hline
        \multicolumn{2}{|l|}{\textbf{Secuencia Base}}                                                                              \\
        \hline
        \multicolumn{2}{|p{15cm}|}{ [1] El usuario accede al panel de administración. }                                            \\
        \hline
        \multicolumn{2}{|p{15cm}|}{ [2] El sistema muestra la lista de Juntas de Cofradías existentes. }                           \\
        \hline
        \multicolumn{2}{|p{15cm}|}{ [3] El usuario selecciona la Junta de Cofradías que desea editar. }                            \\
        \hline
        \multicolumn{2}{|p{15cm}|}{ [4] El usuario modifica los datos de la Junta de Cofradías. }                                  \\
        \hline
        \multicolumn{2}{|p{15cm}|}{ [5] El sistema valida y guarda los cambios. }                                                  \\
        \hline
        \multicolumn{2}{|l|}{\textbf{Secuencia Alternativa}}                                                                       \\
        \hline
        \multicolumn{2}{|p{15cm}|}{ [4.a.1] Cancelación de la operación: El usuario desea cancelar la operación. }                 \\
        \hline
        \multicolumn{2}{|p{15cm}|}{ [5.a.1] Datos incompletos: El sistema muestra los campos requeridos. }                         \\
        \hline
        \textbf{Requisitos Relacionados}                                                                                           & RF-34.                                                                                                                \\
        \hline
    \end{tabular}
    \caption{Caso de uso: Editar Junta de Cofradías}
\end{table}


\begin{table}[h!]
    \centering
    \renewcommand{\arraystretch}{1.5}

    \begin{tabular}{|p{4cm}|p{11cm}|}
        \hline
        \multicolumn{2}{|l|}{\textbf{Caso de uso: Dar de baja Junta de Cofradías}}                                                  \\
        \hline
        \textbf{Resumen}                                                                                                            & El administrador retira una Junta de Cofradías del sistema.                                                          \\
        \hline
        \textbf{Actor}                                                                                                              & Administrador                                                                                                        \\
        \hline
        \textbf{Precondición}                                                                                                       & El usuario tiene rol de Administrador y existe al menos una Junta de Cofradías registrada.                           \\
        \hline
        \textbf{Postcondición}                                                                                                      & La Junta de Cofradías deja de estar disponible en el sistema y la ciudad asociada queda sin Junta asignada.          \\
        \hline
        \multicolumn{2}{|l|}{\textbf{Secuencia Base}}                                                                              \\
        \hline
        \multicolumn{2}{|p{15cm}|}{ [1] El usuario accede al panel de administración. }                                            \\
        \hline
        \multicolumn{2}{|p{15cm}|}{ [2] El sistema muestra la lista de Juntas de Cofradías existentes. }                           \\
        \hline
        \multicolumn{2}{|p{15cm}|}{ [3] El usuario selecciona la Junta de Cofradías que desea eliminar. }                          \\
        \hline
        \multicolumn{2}{|p{15cm}|}{ [4] El sistema comprueba si la Junta de Cofradías tiene cofradías asociadas. }                 \\
        \hline
        \multicolumn{2}{|p{15cm}|}{ [5] El usuario confirma la eliminación. }                                                      \\
        \hline
        \multicolumn{2}{|p{15cm}|}{ [6] El sistema elimina la Junta de Cofradías. }                                                \\
        \hline
        \multicolumn{2}{|l|}{\textbf{Secuencia Alternativa}}                                                                       \\
        \hline
        \multicolumn{2}{|p{15cm}|}{ [4.a.1] Junta con cofradías asociadas: El sistema advierte de que la Junta de Cofradías tiene cofradías registradas antes de continuar. } \\
        \hline
        \multicolumn{2}{|p{15cm}|}{ [5.a.1] Cancelación de la operación: El usuario desea cancelar la eliminación. }               \\
        \hline
        \textbf{Requisitos Relacionados}                                                                                           & RF-34.                                                                                                                \\
        \hline
    \end{tabular}
    \caption{Caso de uso: Dar de baja Junta de Cofradías}
\end{table}

\begin{table}[h!]
    \centering
    \renewcommand{\arraystretch}{1.5}

    \begin{tabular}{|p{4cm}|p{11cm}|}
        \hline
        \multicolumn{2}{|l|}{\textbf{Caso de uso: Editar ciudad}}                                                           \\
        \hline
        \textbf{Resumen}                                                                                                    & El administrador modifica los datos de una ciudad existente.                                                         \\
        \hline
        \textbf{Actor}                                                                                                      & Administrador                                                                                                         \\
        \hline
        \textbf{Precondición}                                                                                               & El usuario tiene rol de Administrador y existe al menos una ciudad registrada.                                       \\
        \hline
        \textbf{Postcondición}                                                                                              & Los datos de la ciudad quedan actualizados en el sistema.                                                            \\
        \hline
        \multicolumn{2}{|l|}{\textbf{Secuencia Base}}                                                                       \\
        \hline
        \multicolumn{2}{|p{15cm}|}{ [1] El usuario accede al panel de administración. }                                     \\
        \hline
        \multicolumn{2}{|p{15cm}|}{ [2] El sistema muestra la lista de ciudades existentes. }                               \\
        \hline
        \multicolumn{2}{|p{15cm}|}{ [3] El usuario selecciona la ciudad que desea editar. }                                 \\
        \hline
        \multicolumn{2}{|p{15cm}|}{ [4] El usuario modifica los datos de la ciudad. }                                       \\
        \hline
        \multicolumn{2}{|p{15cm}|}{ [5] El sistema valida y guarda los cambios. }                                           \\
        \hline
        \multicolumn{2}{|l|}{\textbf{Secuencia Alternativa}}                                                                \\
        \hline
        \multicolumn{2}{|p{15cm}|}{ [4.a.1] Cancelación de la operación: El usuario desea cancelar la operación. }          \\
        \hline
        \multicolumn{2}{|p{15cm}|}{ [5.a.1] Datos incompletos: El sistema muestra los campos requeridos. }                  \\
        \hline
        \textbf{Requisitos Relacionados}                                                                                    & RF-26.                                                                                                                \\
        \hline
    \end{tabular}
    \caption{Caso de uso: Editar ciudad}
\end{table}


\begin{table}[h!]
    \centering
    \renewcommand{\arraystretch}{1.5}

    \begin{tabular}{|p{4cm}|p{11cm}|}
        \hline
        \multicolumn{2}{|l|}{\textbf{Caso de uso: Dar de baja ciudad}}                                                      \\
        \hline
        \textbf{Resumen}                                                                                                    & El administrador retira una ciudad disponible de la aplicación.                                                      \\
        \hline
        \textbf{Actor}                                                                                                      & Administrador                                                                                                         \\
        \hline
        \textbf{Precondición}                                                                                               & El usuario tiene rol de Administrador y existe al menos una ciudad registrada.                                       \\
        \hline
        \textbf{Postcondición}                                                                                              & La ciudad deja de estar disponible en el sistema.                                                                    \\
        \hline
        \multicolumn{2}{|l|}{\textbf{Secuencia Base}}                                                                       \\
        \hline
        \multicolumn{2}{|p{15cm}|}{ [1] El usuario accede al panel de administración. }                                     \\
        \hline
        \multicolumn{2}{|p{15cm}|}{ [2] El sistema muestra la lista de ciudades existentes. }                               \\
        \hline
        \multicolumn{2}{|p{15cm}|}{ [3] El usuario selecciona la ciudad que desea eliminar. }                               \\
        \hline
        \multicolumn{2}{|p{15cm}|}{ [4] El sistema comprueba si la ciudad tiene datos asociados (cofradías, eventos o procesiones). }  \\
        \hline
        \multicolumn{2}{|p{15cm}|}{ [5] El usuario confirma la eliminación. }                                               \\
        \hline
        \multicolumn{2}{|p{15cm}|}{ [6] El sistema elimina la ciudad. }                                                     \\
        \hline
        \multicolumn{2}{|l|}{\textbf{Secuencia Alternativa}}                                                                \\
        \hline
        \multicolumn{2}{|p{15cm}|}{ [4.a.1] Ciudad con datos asociados: El sistema advierte de que la ciudad tiene cofradías, eventos o procesiones registrados antes de continuar. } \\
        \hline
        \multicolumn{2}{|p{15cm}|}{ [5.a.1] Cancelación de la operación: El usuario desea cancelar la eliminación. }        \\
        \hline
        \textbf{Requisitos Relacionados}                                                                                    & RF-26.                                                                                                                \\
        \hline
    \end{tabular}
    \caption{Caso de uso: Dar de baja ciudad}
\end{table}


\begin{table}[h!]
    \centering
    \renewcommand{\arraystretch}{1.5}

    \begin{tabular}{|p{4cm}|p{11cm}|}
        \hline
        \multicolumn{2}{|l|}{\textbf{Caso de uso: Consultar calendario de Semana Santa}}                            \\
        \hline
        \textbf{Resumen}                                                                                           & El ciudadano consulta el calendario con los días de la Semana Santa. \\
        \hline
        \textbf{Actor}                                                                                             & Ciudadano                                                             \\
        \hline
        \textbf{Precondición}                                                                                      & El usuario ha seleccionado una ciudad.                                \\
        \hline
        \textbf{Postcondición}                                                                                     & Se muestra el calendario con los días señalados.                      \\
        \hline
        \multicolumn{2}{|l|}{\textbf{Secuencia Base}}                                                               \\
        \hline
        \multicolumn{2}{|p{15cm}|}{ [1] El usuario accede al calendario desde el menú principal. }                  \\
        \hline
        \multicolumn{2}{|p{15cm}|}{ [2] El sistema muestra el calendario, indicando los días con eventos o procesiones programadas. } \\
        \hline
        \multicolumn{2}{|l|}{\textbf{Secuencia Alternativa}}                                                        \\
        \hline
        \multicolumn{2}{|p{15cm}|}{ [2.a.1] Sin conexión: El sistema muestra la última información almacenada en local. } \\
        \hline
        \textbf{Requisitos Relacionados}                                                                           & RF-35.                                                                \\
        \hline
    \end{tabular}
    \caption{Caso de uso: Consultar calendario de Semana Santa}
\end{table}


\begin{table}[h!]
    \centering
    \renewcommand{\arraystretch}{1.5}

    \begin{tabular}{|p{4cm}|p{11cm}|}
        \hline
        \multicolumn{2}{|l|}{\textbf{Caso de uso: Consultar procesiones de un día}}                                 \\
        \hline
        \textbf{Resumen}                                                                                           & El ciudadano consulta el listado de procesiones programadas en un día concreto. \\
        \hline
        \textbf{Actor}                                                                                             & Ciudadano                                                             \\
        \hline
        \textbf{Precondición}                                                                                      & El usuario está consultando el calendario de Semana Santa.            \\
        \hline
        \textbf{Postcondición}                                                                                     & Se muestra el listado de procesiones de ese día.                      \\
        \hline
        \multicolumn{2}{|l|}{\textbf{Secuencia Base}}                                                               \\
        \hline
        \multicolumn{2}{|p{15cm}|}{ [1] El usuario selecciona un día en el calendario. }                            \\
        \hline
        \multicolumn{2}{|p{15cm}|}{ [2] El sistema muestra el listado de procesiones programadas ese día. }         \\
        \hline
        \multicolumn{2}{|l|}{\textbf{Secuencia Alternativa}}                                                        \\
        \hline
        \multicolumn{2}{|p{15cm}|}{ [2.a.1] Sin procesiones: El sistema indica que no hay procesiones programadas ese día. } \\
        \hline
        \textbf{Requisitos Relacionados}                                                                           & RF-36.                                                                \\
        \hline
    \end{tabular}
    \caption{Caso de uso: Consultar procesiones de un día}
\end{table}


\begin{table}[h!]
    \centering
    \renewcommand{\arraystretch}{1.5}

    \begin{tabular}{|p{4cm}|p{11cm}|}
        \hline
        \multicolumn{2}{|l|}{\textbf{Caso de uso: Acceder a listados desde el menú principal}}                      \\
        \hline
        \textbf{Resumen}                                                                                           & El ciudadano accede a los listados completos de procesiones, pasos, eventos o cofradías. \\
        \hline
        \textbf{Actor}                                                                                             & Ciudadano                                                             \\
        \hline
        \textbf{Precondición}                                                                                      & El usuario ha seleccionado una ciudad.                                \\
        \hline
        \textbf{Postcondición}                                                                                     & Se muestra el listado completo solicitado.                            \\
        \hline
        \multicolumn{2}{|l|}{\textbf{Secuencia Base}}                                                               \\
        \hline
        \multicolumn{2}{|p{15cm}|}{ [1] El usuario accede al menú principal. }                                      \\
        \hline
        \multicolumn{2}{|p{15cm}|}{ [2] El usuario selecciona el tipo de listado (procesiones, pasos, eventos o cofradías). } \\
        \hline
        \multicolumn{2}{|p{15cm}|}{ [3] El sistema muestra el listado completo correspondiente. }                   \\
        \hline
        \multicolumn{2}{|l|}{\textbf{Secuencia Alternativa}}                                                        \\
        \hline
        \multicolumn{2}{|p{15cm}|}{ [3.a.1] Sin elementos disponibles: El sistema muestra un mensaje indicando que no hay elementos de ese tipo. } \\
        \hline
        \textbf{Requisitos Relacionados}                                                                           & RF-37.                                                                \\
        \hline
    \end{tabular}
    \caption{Caso de uso: Acceder a listados desde el menú principal}
\end{table}


\begin{table}[h!]
    \centering
    \renewcommand{\arraystretch}{1.5}

    \begin{tabular}{|p{4cm}|p{11cm}|}
        \hline
        \multicolumn{2}{|l|}{\textbf{Caso de uso: Gestionar perfil de la Junta de Cofradías}}                       \\
        \hline
        \textbf{Resumen}                                                                                           & La Junta de Cofradías consulta o modifica la información de su propio perfil. \\
        \hline
        \textbf{Actor}                                                                                             & Junta de Cofradías                                                    \\
        \hline
        \textbf{Precondición}                                                                                      & El usuario tiene rol de Junta de Cofradías.                           \\
        \hline
        \textbf{Postcondición}                                                                                     & Los datos del perfil quedan actualizados en el sistema.               \\
        \hline
        \multicolumn{2}{|l|}{\textbf{Secuencia Base}}                                                               \\
        \hline
        \multicolumn{2}{|p{15cm}|}{ [1] El usuario accede a la sección de perfil. }                                 \\
        \hline
        \multicolumn{2}{|p{15cm}|}{ [2] El sistema muestra los datos actuales del perfil. }                         \\
        \hline
        \multicolumn{2}{|p{15cm}|}{ [3] El usuario modifica los datos. }                                            \\
        \hline
        \multicolumn{2}{|p{15cm}|}{ [4] El sistema valida y guarda los cambios. }                                   \\
        \hline
        \multicolumn{2}{|l|}{\textbf{Secuencia Alternativa}}                                                        \\
        \hline
        \multicolumn{2}{|p{15cm}|}{ [3.a.1] Cancelación de la operación: El usuario desea cancelar la operación. }  \\
        \hline
        \multicolumn{2}{|p{15cm}|}{ [4.a.1] Datos inválidos: El sistema muestra los errores de validación. }        \\
        \hline
        \textbf{Requisitos Relacionados}                                                                           & RF-39.                                                                \\
        \hline
    \end{tabular}
    \caption{Caso de uso: Gestionar perfil de la Junta de Cofradías}
\end{table}


\begin{table}[h!]
    \centering
    \renewcommand{\arraystretch}{1.5}

    \begin{tabular}{|p{4cm}|p{11cm}|}
        \hline
        \multicolumn{2}{|l|}{\textbf{Caso de uso: Gestionar miembros de las Juntas de Cofradías}}                   \\
        \hline
        \textbf{Resumen}                                                                                           & El administrador da de alta, edita o da de baja a los miembros de una Junta de Cofradías. \\
        \hline
        \textbf{Actor}                                                                                             & Administrador                                                         \\
        \hline
        \textbf{Precondición}                                                                                      & El usuario tiene rol de Administrador y existe al menos una Junta de Cofradías registrada. \\
        \hline
        \textbf{Postcondición}                                                                                     & Los miembros de la Junta de Cofradías quedan actualizados en el sistema. \\
        \hline
        \multicolumn{2}{|l|}{\textbf{Secuencia Base}}                                                               \\
        \hline
        \multicolumn{2}{|p{15cm}|}{ [1] El usuario accede al panel de administración. }                             \\
        \hline
        \multicolumn{2}{|p{15cm}|}{ [2] El usuario selecciona la Junta de Cofradías. }                              \\
        \hline
        \multicolumn{2}{|p{15cm}|}{ [3] El sistema muestra la lista de miembros de esa Junta. }                     \\
        \hline
        \multicolumn{2}{|p{15cm}|}{ [4] El usuario da de alta, edita o da de baja a un miembro. }                   \\
        \hline
        \multicolumn{2}{|p{15cm}|}{ [5] El sistema valida y guarda los cambios. }                                   \\
        \hline
        \multicolumn{2}{|l|}{\textbf{Secuencia Alternativa}}                                                        \\
        \hline
        \multicolumn{2}{|p{15cm}|}{ [4.a.1] Cancelación de la operación: El usuario desea cancelar la operación. }  \\
        \hline
        \multicolumn{2}{|p{15cm}|}{ [5.a.1] Datos inválidos: El sistema muestra los errores de validación. }        \\
        \hline
        \textbf{Requisitos Relacionados}                                                                           & RF-40.                                                                \\
        \hline
    \end{tabular}
    \caption{Caso de uso: Gestionar miembros de las Juntas de Cofradías}
\end{table}